% !TEX root = main.tex
\section{The basic formula of Collatz Sequence}\label{sec:the-basic-formula-of-collatz-sequence}

With the definition of Collatz complement, we can easily find that the equation in \cref{equation:collatz-inductioon-3-2-1} also hold with all general Collatz sequences:
\begin{theorem}~\label{theorem:inductive-formula}
$\forall a \in \widehat{\mathbb{N}}_{{\beta}}$, $\forall k \in \mathbb{N}$ we have
\[ a_k \cdot \beta^{\mathcal{A}_k} = a \cdot\alpha^k + \gamma\cdot\mathfrak{A}_{k}. \]
\end{theorem}
% \begin{proof} Trivial. \end{proof}
\begin{proof}
    First, we have $a_0 \cdot \beta^{\mathcal{A}_0} = a = a \cdot \alpha^{0} + \gamma \cdot \mathfrak{A}_0$ and $a_1 \cdot \beta^{\mathcal{A}_1} = \alpha \cdot a + \gamma = a \cdot \alpha^{1} + \gamma \cdot \mathfrak{A}_1$.\newline
    Second, if $a_k \cdot \beta^{\mathcal{A}_k} = a \cdot\alpha^k + \gamma\cdot\mathfrak{A}_{k}$, then
    \begin{align*}
        a_{k+1} \cdot \beta^{\mathcal{A}_{k+1}}
        =&\ a_{k+1} \cdot \beta^{\val(a;k+1)} \cdot \beta^{\mathcal{A}_{k}} \\
        =&\ \alpha \cdot a_k \cdot \beta^{\mathcal{A}_k} + \gamma \cdot \beta^{\mathcal{A}_k} \\
        % =&\ \alpha \cdot (a \cdot\alpha^k + \gamma\cdot\mathfrak{A}_{k}) + \gamma \cdot \beta^{\mathcal{A}_k} \\
        =&\ a \cdot \alpha^{k+1} + \gamma \cdot \alpha \cdot \mathfrak{A}_{k} + \gamma \cdot \beta^{\mathcal{A}_k} \\
        % =&\ a \cdot \alpha^{k+1} + \gamma \cdot (\alpha \cdot \mathfrak{A}_{k} + \cdot \beta^{\mathcal{A}_k}) \\
        =&\ a \cdot \alpha^{k+1} + \gamma \cdot \mathfrak{A}_{k+1}.
    \end{align*}
    Sum up, we proved this theorem.
\end{proof}
It is very important and fundamental, and we will see later that all types of Collatz sequences are actually some kind of variant expression of the \cref{theorem:inductive-formula}. And this formula also explains why we call $\mathfrak{A}_k$ as the Collatz complement.
\begin{corollary}
$\forall m, n \ge 0$ we have
\[ a_{n+m} \cdot \beta^{\mathcal{A}_{n+m}^{m}} = a_{m} \cdot\alpha^{n} + \gamma\cdot\mathfrak{A}_{n+m}^{m}. \]
\end{corollary}
\begin{corollary}~\label{corollary:inductive-formula-with-reduced}
    $a_{n+m} \cdot \beta^{\mathcal{A}_{n+m}^{m}} = \alpha^{n}\cdot\left(a_m + \frac{\gamma}{\alpha}\cdot\alpha^{m}\cdot\mathfrak{a}_{n+m}^{m}\right)$.
\end{corollary}
\begin{corollary}
Given $a\in\widehat{\mathbb{N}}_\beta$, $\mathfrak{A}$ is of $a$, we have $\beta \nmid \mathfrak{A}_n^m$.
\end{corollary}
\begin{corollary}~\label{corollary:congerunce-with-collatz-complement}
If $\mathcal{A}_k > 0$, then $a \equiv - \frac{\gamma}{\alpha}\cdot\mathfrak{a}_k = - \frac{\gamma}{\alpha^k}\cdot\mathfrak{A}_k \pmod{\beta^{\mathcal{A}_k}}$.
\end{corollary}
% \begin{proof} Trivial with \cref{corollary:inductive-formula-with-reduced}. \end{proof}
\begin{corollary}[Co-uniqueness]~\label{cor:co-uniqueness}
    If there are two numbers $a, b \in \widehat{\mathbb{N}}_{{\beta}}$ such that $\exists n \ge 0$, $\forall m \le n$, there is $\mathcal{A}_m = \mathcal{B}_m$, then $\beta^{\mathcal{A}_n} \mid (a - b)$ and $\alpha^{n} \mid (a_n - b_n)$.
\end{corollary}
% \begin{proof} Trivial with \cref{theorem:inductive-formula}. \end{proof}
\begin{proof}
    Since $\mathcal{A}_m = \mathcal{B}_m$, we have $\mathfrak{A}_{m} = \mathfrak{B}_{m}$.
    With the \cref{theorem:inductive-formula}, we have
    \[ (a_{m} - b_{m}) \cdot \beta^{\mathcal{A}_m} = (a - b) \cdot \alpha^{m}. \]
    Since $\gcd(\alpha, \beta) = 1$, then $\gcd(\beta^{\mathcal{A}_m}, \alpha^{m}) = 1$.
    So we have
    \[\beta^{\mathcal{A}_m} \mid (a - b), \quad \alpha^{m} \mid (a_m - b_m) \]
    holds for all $m\le n$.
\end{proof}
\begin{corollary}
    If $\forall m \ge 0$, there is $\mathcal{A}_m = \mathcal{B}_m$, then we have $a = b$.
\end{corollary}
\begin{corollary}
    If $a = t \cdot p^n + s$, that $t, s\in \mathbb{N}_{\ge1} - \beta\mathbb{N}_{\ge1}, p \in \mathbb{N}_{>1}$, then
    \[a_k \cdot \beta^{\mathcal{A}_k} - \gamma\cdot\mathfrak{A}_k = (t_k \cdot \beta^{\mathcal{T}_k} - \gamma\cdot\mathfrak{T}_k)\cdot p^{n} + (s_k \cdot \beta^{\mathcal{S}_k} - \gamma\cdot\mathfrak{S}_k).\]
\end{corollary}
\begin{corollary}
    If $a = t \cdot \beta^n + s$, that $t, s \in \widehat{\mathbb{N}}_{\beta;\ge1}$ and $n \in \mathbb{N}_{\ge1}$, then
    \begin{align*}
        \mathcal{A}_k &= \mathcal{S}_k, \\
        \mathfrak{A}_k &= \mathfrak{S}_k,\\
        a_k &= t \cdot \alpha^{k} \cdot \beta^{n-\mathcal{S}_k} + s_k
    \end{align*}
    holds during $S_k < n$.
\end{corollary}

It seems like we can introduce the Diophantine equation of $\mathfrak{A}_k$ now, but before that, we have to make sure the $\mathcal{A}_k$ will not finally increased by $0$, i.e.\ fall into those numbers cannot be reduced.
Though we have a clear evidence that $(\alpha, \beta, \gamma) = (3, 2, 1)$ guarantee that there exists at least one triple such that $\mathcal{A}_k \ge n$, we can't say all the triples required in \cref{definition:collatz-sequence} retain this property.
\begin{definition}
Given $m,n \in \mathbb{N}$, when $m \ge 1$, define
\[ [m\mid n] = \left\{a \in \mathbb{N}\ \left|\ \left(a\cdot\alpha^n + \gamma\cdot\sum_{i = 0}^{n-1}\alpha^i\right) \equiv 0 \pmod{\beta^{m}}\right.\right\}. \]
Specially, when $m = 0$, define $[0\mid n] = \mathbb{N}$.
\end{definition}
\begin{proposition}~\label{proposition:equivalence-of-classes}
    $a \in [m\mid n] \Longleftrightarrow a\equiv -\frac{\gamma}{\alpha}\cdot\sum_{i=0}^{n-1}\alpha^{-i}\pmod{\beta^{m}}$.
\end{proposition}
\begin{proposition}
    $ l \ge m \Longleftrightarrow [l\mid n] \subset [m\mid n]$.
\end{proposition}
\begin{proposition}~\label{proposition:avaliable-in-finity-step}%
    $\exists P \in \mathbb{N}_{\ge1}$, such that $[m \mid n + P ] = [m\mid n]$.
\end{proposition}
\begin{proof}
    Let $P = \ord_{\beta^{m}}(\alpha) \cdot \prod_{p\in\Pi(\mathbb{P} \mid \beta)}\raf_{p}(\alpha-1)$.

    First, since $\gcd(\alpha, \beta) = 1$, for those $p \in \mathbb{P}$ that $p \mid \beta$, we have
    \begin{align*}
        \fac_{p}(\alpha^P - 1) & \ge \fac_{p}\left(\prod_{p\in\Pi(\mathbb{P} \mid \beta)}\raf_{p}(\alpha-1)\right) + \fac_{p}(\alpha^{\ord_{\beta^{m}}(\alpha)} - 1) \\
                                 & \ge \fac_{p}(\alpha-1) + m \cdot \fac_{p}\beta,
    \end{align*}
    Then we have
    \begin{align*}
        \fac_{p}\left(\sum_{i=0}^{P - 1}\alpha^{-i}\right) & = \fac_{p}\left(\frac{\alpha(\alpha^{P} - 1)}{\alpha^{P}(\alpha - 1)}\right) \\
                &= \fac_{p}(\alpha^P - 1) - \fac_{p}(\alpha-1) \\
                &\ge m \cdot \fac_{p}\beta.
    \end{align*}
    That means $(\raf_{p_i}\beta)^{m} \left\vert\ \sum_{i=0}^{P - 1}\alpha^{-i} \right. $.
    Thus, we have $\beta^{m} \left\vert\ \sum_{i=0}^{P - 1}\alpha^{-i} \right. $.
    Here, $\alpha^{-1}$ is defined in terms of the modulus of $\beta^{m}$. \newline
    Now suppose $a \in [m \mid n + P]$.
    From \cref{proposition:equivalence-of-classes}, we have
    \begin{align*}
        a   &\equiv -\frac{\gamma}{\alpha}\cdot\sum_{i=0}^{P+n-1}\alpha^{-i}    &&\pmod{\beta^{m}} \\
            &\equiv -\frac{\gamma}{\alpha}\cdot\sum_{i=0}^{P - 1}\alpha^{-i} - \frac{\gamma}{\alpha}\cdot\sum_{i=P}^{P + n - 1}\alpha^{-i}    &&\pmod{\beta^{m}} \\
            &\equiv - \frac{\gamma}{\alpha}\cdot\alpha^{-P} \cdot\sum_{i=0}^{n - 1}\alpha^{-i}    &&\pmod{\beta^{m}} \\
            &\equiv -\frac{\gamma}{\alpha}\cdot\sum_{i=0}^{n - 1}\alpha^{-i}    &&\pmod{\beta^{m}}
    \end{align*}
    which means $a \in [ m \mid n + P] \Longleftrightarrow a \in [m \mid n]$.
    But this $P$ is not always the minimal period.
\end{proof}
\begin{definition}
    Denote $[m] = \bigcup_{n=0}^{\infty}[m\mid n]$ and $\bar{[m]} = \mathbb{N} - [m]$.
\end{definition}
\begin{proposition}
    $[0] = \mathbb{N}$ and $\bar{[0]} = \varnothing$.
\end{proposition}
\begin{proposition}
    $a \in \bar{[1]} \Longrightarrow \alpha \cdot a + \gamma \in \bar{[1]}$.
    On the contrary, if $a\equiv \gamma \pmod{\alpha}$, then $a \in \bar{[1]} \Longrightarrow \frac{a - \gamma}{\alpha} \in \bar{[1]}$.
\end{proposition}
\begin{theorem}~\label{theorem:requirements-of-avaliablity}
    $\lim_{k\to\infty} \mathcal{A}_k = \infty\Longleftrightarrow\forall k\in \mathbb{N}, a_k \in [1]$.
\end{theorem}
\begin{proof}
\begin{itemize}
    \item[$\Longrightarrow$)]
        If $\lim_{k\to\infty} \mathcal{A}_k = \infty$, then $\forall k \in \mathbb{N}$, $\exists n_{k} \in \mathbb{N}_{\ge1}$, $\val(a;k + n_{k}) > 0$,
        that means $a_{k} \in [\val(a;k + n_{k}) | n_{k}] \subset [1]$.
    \item[$\Longleftarrow$)] 
        If $\forall k \in \mathbb{N}$, there is $a_k \in [1]$, then $\exists n_{k} \in \mathbb{N}_{\ge1}$ and from \cref{proposition:avaliable-in-finity-step} we can suppose $n_k \le P$,
        such that $a_k \in [1 \mid n_{k}]$ for all $k \in \mathbb{N}$.
        Thus, we have $\mathcal{A}_k \ge \frac{k}{P} \to \infty\ (k \to \infty)$.
\end{itemize}
Sum up, we proved this theorem.
\end{proof}
\begin{corollary}
    If $\lim_{k\to\infty}\mathcal{A}_k = \infty$, then $\frac{\mathcal{A}_k}{k} \ge \frac{1}{P}$ when $k\to\infty$.
\end{corollary}
\begin{conjecture}
    $\exists c(\beta) \ge 0$, $\mu\{a \in\widehat{\mathbb{Z}}_{{\beta}} \mid \lim_{k\to\infty}\frac{\mathcal{A}_k}{k} \ge c \}  = 0$.
\end{conjecture}

\cref{theorem:requirements-of-avaliablity} present that, for general triplet arguments, not all initial numbers can be reduced infinitely many times.
But these propositions above provide a basic way to determine whether a number can be reduced, especially since we can easily filter out many initial numbers that cannot. We have
\begin{example}
    Suppose $(\alpha, \beta, \gamma) = (3, 2, 1)$.
    $\forall a \in \widehat{\mathbb{N}}_{2}$, $a\in[1]$, and $1 \le \frac{\mathcal{A}_k}{k} \le 2$ thus $\lim_{k\to\infty}\mathcal{A}_k = \infty$.
\end{example}

However, even if $A_{k}\to \infty$, we can't say that $a_k$ will eventually fall into a cycle.
For example, when $\alpha = 11$, $\beta = 2$ and $\gamma = 1$, for any number, the sequence may grow exponentially as $k$ increases, while $\alpha = 3$ will decrease to 1, at least for the numbers we observe.
The reason for this is that the distribution of reduce always only relates to $\beta$, and its change with $\alpha$ is not significant.
But the growth of sequence numbers accelerates as $\alpha$ increases. $11$ is too large, $\beta=2$ and $\gamma=1$ are dwarfed with it.

Therefore, we can only find these cycle numbers by letting $a_{t} = a$.
According to \cref{theorem:inductive-formula}, that is
\[ a = \frac{\gamma \cdot \mathfrak{A}_{t}}{\beta^{\mathcal{A}_{t}} - \alpha^{t}}. \]
For this equation, there probably are some integer solutions from \cref{proposition:trivial-periodic-sequence}, but not all integer solution of it follows that proposition.
For example, $a=13$ has a 4-tempus cycle under $\alpha=3$, $\beta=5$, $\gamma=1$, that $\val(13;1) = 1$, $\val(13;2) = 2$, $\val(13;3) = 0$ and $\val(13;4) = 0$.
Obviously, $\val(13;i)$ are not same.
Therefore, we have to explore the properties of this fractional expression independently.
